\RequirePackage{pdfmanagement}
\SetKeys[document/metadata]{pdfversion=1.7, pdfstandard=A-2u, lang=en-GB}
\input{preamble/pdfA-compliance/glyphtounicode}
\documentclass[11pt,a4paper]{report}

\newcommand*{\PaperTitle}{No limit cycle in the Banaji--Boros--Hofbauer three-reaction example}
\newcommand*{\PaperTitlePlain}{No limit cycle in the Banaji-Boros-Hofbauer three-reaction example}
\newcommand*{\PaperShortTitle}{No limit cycle in the BBH example}
\newcommand*{\PaperSubtitle}{chemistry-open-problems \textperiodcentered{} 18 October 2026 \textperiodcentered{} preliminary}
\newcommand*{\PaperCredits}{Dulac argument and subfamily: Claude Opus 5.5 exploration agents (scout and verifier, effort high), 17 Oct 2026 \textperiodcentered{} Certificate and manuscript: Claude Opus 5.5 (effort xhigh) \textperiodcentered{} Independent verification: pending}
\newcommand*{\PaperKeywords}{mass-action kinetics, limit cycles, Bendixson--Dulac criterion, Andronov--Hopf bifurcation}
\input{chemop-preamble}
\addbibresource{references.bib}

\begin{document}
\makepapertitle

\begin{paperabstract}
Banaji, Boros and Hofbauer classified the three-reaction planar quadratic mass-action networks that admit an
Andronov--Hopf bifurcation, and asked whether the networks of their cases 9 and 10 without such a bifurcation can
nevertheless have a limit cycle; as an example they asked about the network
\(2X\to3X+2Y,\ X+Y\to0,\ Y\to X+Y\). We show that this system has no periodic orbit for any rate constants, using a
Dulac function that turns the divergence into a negative definite quadratic form. The same argument shows that in
the subfamily \(2X\to(2+c_1)X+d_1Y,\ X+Y\to0,\ Y\to(1+c_3)X+Y\), periodic orbits can only occur in the
Andronov--Hopf region \(c_1<d_1<2c_1\).
\end{paperabstract}

\section*{The question}

Banaji, Boros and Hofbauer~\autocite{BBH2024} classify the three-reaction, planar, quadratic mass-action systems that
admit an Andronov--Hopf bifurcation (their Theorem~9, cases~9 and~10). In their Conclusions (§6) they write:

\begin{openquote}
``However, we leave it open whether the networks in cases 9 and 10 admitting no Andronov--Hopf bifurcations can
have a limit cycle. For example, we do not know whether the system \(2X\to3X+2Y\) (\(\kappa_1\)), \(X+Y\to0\)
(\(\kappa_2\)), \(Y\to X+Y\) (\(\kappa_3\)), \(\dot x=\kappa_1x^2-\kappa_2xy+\kappa_3y\),
\(\dot y=2\kappa_1x^2-\kappa_2xy\) admits a limit cycle.''
\end{openquote}

We settle the example (it has no periodic orbit at all, for any rate constants) and, with the same argument, the
subfamily of case-9 networks to which it belongs.

\begin{theorem}\label{thm:example}
For all \(\kappa_1,\kappa_2,\kappa_3>0\), the system
\[
\dot x=\kappa_1x^2-\kappa_2xy+\kappa_3y,\qquad \dot y=2\kappa_1x^2-\kappa_2xy
\]
has no periodic orbit in the closed positive quadrant. In particular it has no limit cycle.
\end{theorem}

\begin{theorem}\label{thm:family}
Let \(c_1,c_3,d_1\) be positive integers and consider the networks
\[
2X\xrightarrow{\kappa_1}(2+c_1)X+d_1Y,\qquad X+Y\xrightarrow{\kappa_2}0,\qquad Y\xrightarrow{\kappa_3}(1+c_3)X+Y,
\]
i.e.\ case~9 of~\autocite[Theorem~9]{BBH2024} with reaction vectors \(c=(c_1,-1,c_3)\), \(d=(d_1,-1,0)\). The
mass-action system
\[
\dot x=c_1\kappa_1x^2-\kappa_2xy+c_3\kappa_3y,\qquad \dot y=d_1\kappa_1x^2-\kappa_2xy
\]
has a periodic orbit for some \(\kappa\) only if \(c_1<d_1<2c_1\), which is exactly the region where it undergoes an
Andronov--Hopf bifurcation~\autocite[Theorem~9]{BBH2024}. For \(d_1\le c_1\) or \(d_1\ge 2c_1\) it has no periodic
orbit for any \(\kappa>0\).
\end{theorem}

The example of \zcref{thm:example} is the case \(c_1=c_3=1\), \(d_1=2\), on the boundary \(d_1=2c_1\) of the Hopf
region.

\section{Proofs}

\begin{proof}[of \zcref{thm:family}, case \(d_1\ge2c_1\)]
Write \(F=(f,g)\) for the vector field and take the Dulac function \(B(x,y)=y^{-b}\) on \(\mathbb R^2_{>0}\), with
\[
b=1+\frac{2\kappa_1(d_1-c_1)}{\kappa_2}.
\]
A direct computation gives
\[
y^{\,b+1}\,\operatorname{div}(BF)=Q(x,y):=\frac{d_1\kappa_1\bigl(2(c_1-d_1)\kappa_1-\kappa_2\bigr)}{\kappa_2}\,x^2
+2d_1\kappa_1\,xy-\kappa_2\,y^2 ,
\]
a quadratic form with \(y^2\) coefficient \(-\kappa_2<0\) and discriminant
\[
\Delta=4d_1\kappa_1\bigl(\kappa_1(2c_1-d_1)-\kappa_2\bigr).
\]
(The exponent \(b\) minimises \(\Delta\) among Dulac functions \(y^{-b}\).) If \(d_1\ge 2c_1\), then
\(\kappa_1(2c_1-d_1)-\kappa_2\le-\kappa_2<0\), so \(\Delta<0\) and \(Q\) is negative definite: \(\operatorname{div}(BF)<0\)
on the open quadrant. Since \(\mathbb R^2_{>0}\) is simply connected, the Bendixson--Dulac criterion excludes periodic
orbits contained in it.

On \(\{y=0,\ x>0\}\) we have \(\dot y=d_1\kappa_1x^2>0\), on \(\{x=0,\ y>0\}\) we have \(\dot x=c_3\kappa_3y>0\), and the
origin is an equilibrium. A periodic orbit in the closed quadrant cannot contain an equilibrium, and cannot touch a
boundary point where the field points strictly into the open quadrant (it would have to reach that point from the
inside). Hence every periodic orbit would lie in the open quadrant, which is impossible.
\end{proof}

\begin{proof}[of \zcref{thm:family}, case \(d_1\le c_1\)]
The stoichiometric matrix \(\Gamma=\begin{pmatrix}c_1&-1&c_3\\ d_1&-1&0\end{pmatrix}\) has rank~2 (its rows are not
proportional because \(c_3\neq0\)), so \(\ker\Gamma\) is spanned by \(u=c\times d=(c_3,\;c_3d_1,\;d_1-c_1)\). A positive
equilibrium \(x\) requires a positive flux vector \(v(x)\in\ker\Gamma\cap\mathbb R^3_{>0}\), i.e.\ \(u\) or \(-u\)
positive, which fails when \(d_1\le c_1\). So there is no positive equilibrium. A periodic orbit in the open quadrant
would enclose an equilibrium (index theory), which would be positive; hence there is none, and the boundary argument
above excludes orbits touching the axes.
\end{proof}

\begin{proof}[of \zcref{thm:example}]
This is \(c_1=c_3=1\), \(d_1=2\): \(b=1+2\kappa_1/\kappa_2\),
\(Q=-\frac{2\kappa_1(2\kappa_1+\kappa_2)}{\kappa_2}x^2+4\kappa_1xy-\kappa_2y^2\) and \(\Delta=-8\kappa_1\kappa_2<0\).
\end{proof}

\begin{remark}
In the example, the unique positive equilibrium is \((2\kappa_3/\kappa_2,\ 4\kappa_1\kappa_3/\kappa_2^2)\), with
\(\det J=4\kappa_1\kappa_3^2/\kappa_2>0\) and \(\operatorname{tr}J=-2\kappa_3<0\).

For case~9 with \(c_3>0\), \autocite[Theorem~9]{BBH2024} gives an Andronov--Hopf bifurcation iff
\(\tfrac12(d_3/c_3+d_1/c_1)<d_2/c_2<d_1/c_1\). In the subfamily, \(d_2/c_2=1\) and \(d_3=0\), so the condition reads
\(c_1<d_1<2c_1\); for the example, \(\tfrac12(0+2)=1<1<2\) fails by equality. \zcref[S]{thm:family} shows that in this
subfamily ``no Andronov--Hopf bifurcation'' implies ``no periodic orbit''.
\end{remark}

\section{What remains open}

The general question of~\autocite{BBH2024}, for all case-9 and case-10 networks without an Andronov--Hopf
bifurcation (case~9 with \(d_2\neq-1\) or \(d_3\neq0\), including the branch \(c_3=0\), and all of case~10), is not
settled here. The same strategy with Dulac functions \(x^\alpha y^\beta\) reduces it to sign conditions on a
cubic-plus-quadratic form.

\section{Computer verification}

\texttt{scripts/certify\_dulac.py} checks every identity above symbolically (sympy): the Dulac computation, the
coefficients and discriminant in the subfamily and in the example, the vector \(u=c\times d\), the equilibrium and its
Jacobian, and the boundary behaviour.

\paperreferences
\end{document}
